\documentclass[10pt,a4paper]{amsart}
\usepackage[margin=19mm]{geometry}
\usepackage{amsmath,amssymb,mathtools}
\usepackage[scr=rsfs,cal=euler]{mathalfa}
\usepackage{xfrac}
\usepackage[unicode,hidelinks,bookmarks=false]{hyperref}
\newcommand{\ie}{i.e.\ }
\newcommand{\cf}{cf.\ }
\newcommand{\resp}{respectively\ }
\newcommand{\Subsection}{\subsection}
\providecommand{\nobreakdash}{\nobreak\hbox{-}}
\setlength{\emergencystretch}{2em}
\multlinegap0pt
\usepackage{bm}
\usepackage{amssymb}
\usepackage{mathtools}
\usepackage{xcolor}
\usepackage{tikz}
\usepackage[shortlabels]{enumitem}
\newtheorem{prop}{Proposition}
\newcommand{\cC}{\mathcal C}
\newcommand{\id}{\mathbf 1}
\title{Virasoro Constraints for Orbifold Curves}
\author{Jianxun Hu, Pengrong Jin, Hua-Zhong Ke, Jingwei Lu, Hsian-Hua Tseng, Longting Wu}
\date{Source: arXiv:2609.01300v1 (CC BY 4.0)}
\begin{document}
\maketitle

\noindent\emph{Source and licence.} Virasoro Constraints for Orbifold Curves. Version arXiv:2609.01300v1 (CC BY 4.0), distributed by arXiv under the Creative Commons Attribution 4.0 International licence, http://creativecommons.org/licenses/by/4.0/, as recorded in the arXiv licence metadata for this exact version and shown on the abstract page https://arxiv.org/abs/2609.01300v1. Full licence notice: \texttt{licenses/arxiv-2609.01300-CC-BY-4.0.txt}.\par\smallskip

\noindent\textit{Editorial scope.} Counted unit: the article's prop prop:virasoro-defect-degeneration (source0.tex lines 616--689). The article's complete original proof of it is delivered with it (source0.tex lines 691--1071, one \texttt{proof} environment of 381 lines). No mathematics is written, repaired or re-derived: the statement and the proof are the article's own lines, in the article's own notation, and no retained line is substituted or re-flowed.  Retained uncounted as the source-local context the proof needs: lines 500--524.  Nothing was omitted from any retained span, so the package carries no omission record.  No retained line contains a citation, so the wrapper carries no bibliography and no bibliography entry.  No retained line contains a figure environment, an image-inclusion command or drawing code, so no image is delivered and the package declares no asset dependency.  Editorial formatting only, in addition: an AMS \texttt{amsart} preamble that restates the article's own declarations and macros used by the retained lines, namely the environments \texttt{prop} on separate counters, together with the standard packages those lines need.  The retained environments print in this document as Proposition 1 in order of appearance, whereas the article prints the same items with the numbering of its own full document.  The complete native source of the article as distributed in the arXiv e-print for this exact version is bundled unchanged as \texttt{sources/209096/source0.tex}.
\begin{equation}
\label{eq:relative-orbifold-degeneration}
\begin{aligned}
\left\langle
\prod_{i\in N}\tau_{k_i}(\gamma_i)
\ \middle|\
\boldsymbol\eta',\boldsymbol\eta_0
\right\rangle^{\bullet,\,\cC/Q}_{d}
={}&
\sum_{\mu\vdash d}\mathfrak z(\mu)
\sum_{N=N'\sqcup N_0}(-1)^{\epsilon(N',N_0)}
\\[-1mm]
&\hspace{-35mm}\cdot
\left\langle
\prod_{i\in N'}\tau_{k_i}(\gamma_i)
\ \middle|\
\boldsymbol\eta',\mu
\right\rangle^{\bullet,\,\cC'/(Q'\cup\{p\})}_{d}
\left\langle
\prod_{i\in N_0}\tau_{k_i}(\gamma_{i})
\ \middle|\
\boldsymbol\eta_0,\mu
\right\rangle^{\bullet,\,\cC_0/(Q_0\cup\{p\})}_{d}.
\end{aligned}
\end{equation}

\begin{prop}
\label{prop:virasoro-defect-degeneration}
Use the relative-data notation fixed at the beginning of this section.
Assume that the coarse curve underlying \(\cC_0\)
has genus zero.
Let
\[
\mathbf B=
\prod_{i=1}^{u}\tau_{a_i}(\id)
\prod_{j=1}^{v}\tau_{b_j}(\omega)
\mathbf A'\mathbf A_0,
\]
where \(\mathbf A'\) and \(\mathbf A_0\) are the ordered products of the
remaining insertions supported on \(\cC'\) and \(\cC_0\), respectively. We fix a partition \(\{1,\ldots,v\}=T\sqcup T'\), such that for $j\in T$, the representative of $\omega$ in $\tau_{b_j}(\omega)$ is chosen to be $(\omega',\omega_0)=(0,\omega)$, while for $j\in T'$, the representative of $\omega$ in $\tau_{b_j}(\omega)$ is chosen to be $(\omega',\omega_0)=(\omega,0)$. For \(S\subset\{1,\ldots,u\}\), set
\[
\mathbf B'(S)
=
\prod_{i\notin S}\tau_{a_i}(\id)
\prod_{j\notin T}\tau_{b_j}(\omega)
\mathbf A',
\qquad
\mathbf B_0(S)
=
\prod_{i\in S}\tau_{a_i}(\id)
\prod_{j\in T}\tau_{b_j}(\omega)
\mathbf A_0.
\]

Then, for every \(k\geq-1\),
\begin{equation}
\label{eq:virasoro-defect-degeneration}
\begin{aligned}
&
\Delta_{k,d}^{\cC/Q}
\left(
\mathbf B
\ \middle|\
\boldsymbol\eta
\right)
\\
&=
\sum_{S\subset\{1,\ldots,u\}}
\sum_{\mu\vdash d}
\mathfrak z(\mu)
\Bigg[
\Delta_{k,d}^{\cC'/(Q'\cup\{p\})}
\left(
\mathbf B'(S)
\ \middle|\
\boldsymbol\eta',\mu
\right)\cdot
\left\langle
\mathbf B_0(S)
\ \middle|\
\boldsymbol\eta_0,\mu
\right\rangle^{\bullet,\,\cC_0/(Q_0\cup\{p\})}_{d}
\\
&\hspace{25mm}
+
\left\langle
\mathbf B'(S)
\ \middle|\
\boldsymbol\eta',\mu
\right\rangle^{\bullet,\,\cC'/(Q'\cup\{p\})}_{d}\cdot
\Delta_{k,d}^{\cC_0/(Q_0\cup\{p\})}
\left(
\mathbf B_0(S)
\ \middle|\
\boldsymbol\eta_0,\mu
\right)
\Bigg].
\end{aligned}
\end{equation}
\end{prop}

\begin{proof}

We use the following notations in the proof. Recall that an insertion type is of the form \(\tau_\ell(\gamma)\).
Let \(A\) and \(B\) be insertion types.
For an ordered insertion list \(\mathbf D\), let
\(\Lambda_A(\mathbf D)\) be the set of markings with insertion of type
\(A\) in \(\mathbf D\).  For
\(\lambda\in\Lambda_A(\mathbf D)\), let
\(\mathbf D_{\lambda:A\rightsquigarrow B}\) be the new ordered insertion list obtained by replacing the type \(A\) at the
marking \(\lambda\) in $\mathbf D$ by \(B\).  Let
\(\mathbf D\sqcup\{A\}\) be the  new ordered insertion list obtained by adding one new insertion of type \(A\) after $\mathbf D$.
For two distinct markings \(\lambda,\nu\) of \(\mathbf D\), let
\(\mathbf D\setminus\{\lambda,\nu\}\) be the new  ordered insertion list obtained by deleting the corresponding insertions. 

With these conventions, all four identities below hold when \(A\) and \(B\) are
both even.  The second identity also holds when \(A\) and \(B\) are
both odd:
\begin{align*}
\mathbf D!\,[\mathbf t_{\mathbf D}]\frac{\partial}{\partial\mathbf t_A} Z_d^{\cC/Q}[\boldsymbol\eta]&=\left\langle\mathbf D\sqcup\{A\}\ \middle|\ \boldsymbol\eta\right\rangle_d^{\bullet,\,\cC/Q},\\
\mathbf D!\,[\mathbf t_{\mathbf D}]\,\mathbf t_A\frac{\partial}{\partial\mathbf t_B} Z_d^{\cC/Q}[\boldsymbol\eta]&=\sum_{\lambda\in\Lambda_A(\mathbf D)}\left\langle\mathbf D_{\lambda:A\rightsquigarrow B}\ \middle|\ \boldsymbol\eta\right\rangle_d^{\bullet,\,\cC/Q},\\
\mathbf D!\,[\mathbf t_{\mathbf D}]\frac{\partial^2 }{\partial\mathbf t_A\,\partial\mathbf t_B}Z_d^{\cC/Q}[\boldsymbol\eta]&=\left\langle\mathbf D\sqcup\{A,B\}\ \middle|\ \boldsymbol\eta\right\rangle_d^{\bullet,\,\cC/Q},\\
\mathbf D!\,[\mathbf t_{\mathbf D}]\,\mathbf t_A\mathbf t_B Z_d^{\cC/Q}[\boldsymbol\eta]&=\sum_{\substack{\lambda\in\Lambda_A(\mathbf D),\,\nu\in\Lambda_B(\mathbf D)\\ \lambda\ne\nu}}\left\langle\mathbf D\setminus\{\lambda,\nu\}\ \middle|\ \boldsymbol\eta\right\rangle_d^{\bullet,\,\cC/Q}.
\end{align*}

In what follows, we suppress target-independent scalar coefficients of the operators for ease of notation, but retain
the coefficients involving \(\chi_{\log}\) which is target-dependent. Note that
\begin{equation}\label{eq:chi-log-degeneration-in-proof}
\chi_{\log}(\cC,Q)=\chi_{\log}(\cC',Q'\cup\{p\})+\chi_{\log}(\cC_0,Q_0\cup\{p\}).
\end{equation}

For \(\partial/\partial t^0_{k+1}\), we add
\(\tau_{k+1}(\id)\) in the insertions. Now \eqref{eq:relative-orbifold-degeneration} gives
\[
\begin{aligned}
&\left\langle\mathbf B\sqcup\{\tau_{k+1}(\id)\}
 \ \middle|\ \boldsymbol\eta',\boldsymbol\eta_0\right\rangle_d^{\bullet,\,\cC/Q}
={}\sum_{S\subset\{1,\ldots,u\}}\sum_{\mu\vdash d}\mathfrak z(\mu)\Bigg[
\\[-2mm]
&\quad\left\langle\mathbf B'(S)\sqcup\{\tau_{k+1}(\id)\}
 \ \middle|\ \boldsymbol\eta',\mu\right\rangle_d^{\bullet,\,\cC'/(Q'\cup\{p\})}
 \left\langle\mathbf B_0(S)\ \middle|\ \boldsymbol\eta_0,\mu\right\rangle_d^{\bullet,\,\cC_0/(Q_0\cup\{p\})}
\\[-1mm]
&\quad+\left\langle\mathbf B'(S)\ \middle|\ \boldsymbol\eta',\mu\right\rangle_d^{\bullet,\,\cC'/(Q'\cup\{p\})}
 \left\langle\mathbf B_0(S)\sqcup\{\tau_{k+1}(\id)\}
 \ \middle|\ \boldsymbol\eta_0,\mu\right\rangle_d^{\bullet,\,\cC_0/(Q_0\cup\{p\})}\Bigg].
\end{aligned}
\]

For \(-\chi_{\log}(\cC,Q)\,\partial/\partial t^1_k\), we add
\(\tau_k(\omega)\) in the insertions. From \eqref{eq:chi-log-degeneration-in-proof}, we have
\begin{align*}
 &-\chi_{\log}(\cC,Q)
 \left\langle\mathbf B\sqcup\{\tau_k(\omega)\}
 \ \middle|\ \boldsymbol\eta',\boldsymbol\eta_0\right\rangle_d^{\bullet,\,\cC/Q}\\
 &=
 -\chi_{\log}(\cC',Q'\cup\{p\})
 \left\langle\mathbf B\sqcup\{\tau_k(\omega)\}
 \ \middle|\ \boldsymbol\eta',\boldsymbol\eta_0\right\rangle_d^{\bullet,\,\cC/Q}\\
 &\hspace*{14mm}-\chi_{\log}(\cC_0,Q_0\cup\{p\})
 \left\langle\mathbf B\sqcup\{\tau_k(\omega)\}
 \ \middle|\ \boldsymbol\eta',\boldsymbol\eta_0\right\rangle_d^{\bullet,\,\cC/Q}\\
 &=(I)+(II).
\end{align*}
For the added $\tau_k(\omega)$ in $(I)$, we choose the representative $(\omega',\omega_0)=(\omega,0)$, and for the added $\tau_k(\omega)$ in $(II)$, we choose the representative $(\omega',\omega_0)=(0,\omega)$. Now \eqref{eq:relative-orbifold-degeneration} gives
\[
\begin{aligned}
&-\chi_{\log}(\cC,Q)
 \left\langle\mathbf B\sqcup\{\tau_k(\omega)\}
 \ \middle|\ \boldsymbol\eta',\boldsymbol\eta_0\right\rangle_d^{\bullet,\,\cC/Q}
={}\sum_{S\subset\{1,\ldots,u\}}\sum_{\mu\vdash d}\mathfrak z(\mu)\Bigg[
\\[-2mm]
&\quad-\chi_{\log}(\cC',Q'\cup\{p\})
 \left\langle\mathbf B'(S)\sqcup\{\tau_k(\omega)\}
 \,\middle|\,\boldsymbol\eta',\mu\right\rangle_d^{\bullet,\,\cC'/(Q'\cup\{p\})}
 \left\langle\mathbf B_0(S)\,\middle|\,\boldsymbol\eta_0,\mu\right\rangle_d^{\bullet,\,\cC_0/(Q_0\cup\{p\})}
\\[-1mm]
&\quad-\chi_{\log}(\cC_0,Q_0\cup\{p\})
 \left\langle\mathbf B'(S)\,\middle|\,\boldsymbol\eta',\mu\right\rangle_d^{\bullet,\,\cC'/(Q'\cup\{p\})}
 \left\langle\mathbf B_0(S)\sqcup\{\tau_k(\omega)\}
 \,\middle|\,\boldsymbol\eta_0,\mu\right\rangle_d^{\bullet,\,\cC_0/(Q_0\cup\{p\})}\Bigg].
\end{aligned}
\]

For \(t_A\partial/\partial t_B\), where \(A=\tau_\ell(\gamma)\), \(B=\tau_{\ell+k}(\gamma)\) and
\(\gamma\in\{\id,\omega,\alpha_j,\beta_j,\phi_{i,a}\}\), we replace
\(A\) in \(\mathbf B\) by \(B\).  Now \eqref{eq:relative-orbifold-degeneration} gives
\[
\begin{aligned}
&\sum_{\lambda\in\Lambda_A(\mathbf B)}
 \left\langle\mathbf B_{\lambda:A\rightsquigarrow B}
 \ \middle|\ \boldsymbol\eta',\boldsymbol\eta_0\right\rangle_d^{\bullet,\,\cC/Q}
={}\sum_{S\subset\{1,\ldots,u\}}\sum_{\mu\vdash d}\mathfrak z(\mu)\Bigg[
\\[-2mm]
&\quad\sum_{\lambda\in\Lambda_A(\mathbf B'(S))}
 \left\langle\mathbf B'(S)_{\lambda:A\rightsquigarrow B}
 \ \middle|\ \boldsymbol\eta',\mu\right\rangle_d^{\bullet,\,\cC'/(Q'\cup\{p\})}
 \left\langle\mathbf B_0(S)\ \middle|\ \boldsymbol\eta_0,\mu\right\rangle_d^{\bullet,\,\cC_0/(Q_0\cup\{p\})}
\\[-1mm]
&\quad+\left\langle\mathbf B'(S)\ \middle|\ \boldsymbol\eta',\mu\right\rangle_d^{\bullet,\,\cC'/(Q'\cup\{p\})}
 \sum_{\lambda\in\Lambda_A(\mathbf B_0(S))}
 \left\langle\mathbf B_0(S)_{\lambda:A\rightsquigarrow B}
 \ \middle|\ \boldsymbol\eta_0,\mu\right\rangle_d^{\bullet,\,\cC_0/(Q_0\cup\{p\})}\Bigg].
\end{aligned}
\]
% If \(\gamma=\alpha_j\) or \(\beta_j\), its restriction to \(\cC_0\)
% vanishes because \(\cC_0\) has genus zero; hence only the contribution in
% which the changed insertion lies on \(\cC'\) is nonzero.  If
% \(\gamma=\phi_{i,a}\), the class is supported at \(p_i\), so only the
% contribution from the component containing \(p_i\) is nonzero.

For \(\chi_{\log}(\cC,Q)t^0_\ell\partial_{t^1_{\ell+k-1}}\), setting
\(A=\tau_\ell(\id)\), \(B=\tau_{\ell+k-1}(\omega)\), we replace \(A\) in $\mathbf B$ by \(B\).
Following similar approach for \(-\chi_{\log}(\cC,Q)\,\partial/\partial t^1_k\), we get
\[
\begin{aligned}
&\chi_{\log}(\cC,Q)\sum_{\lambda\in\Lambda_A(\mathbf B)}
 \left\langle\mathbf B_{\lambda:A\rightsquigarrow B}
 \ \middle|\ \boldsymbol\eta',\boldsymbol\eta_0\right\rangle_d^{\bullet,\,\cC/Q}
={}\sum_{S\subset\{1,\ldots,u\}}\sum_{\mu\vdash d}\mathfrak z(\mu)\Bigg[
\\[-2mm]
&\quad\chi_{\log}(\cC',Q'\cup\{p\})
 \sum_{\lambda\in\Lambda_A(\mathbf B'(S))}
 \left\langle\mathbf B'(S)_{\lambda:A\rightsquigarrow B}
 \ \middle|\ \boldsymbol\eta',\mu\right\rangle_d^{\bullet,\,\cC'/(Q'\cup\{p\})}
 \left\langle\mathbf B_0(S)\ \middle|\ \boldsymbol\eta_0,\mu\right\rangle_d^{\bullet,\,\cC_0/(Q_0\cup\{p\})}
\\[-2mm]
&+\chi_{\log}(\cC_0,Q_0\cup\{p\})
 \sum_{\lambda\in\Lambda_A(\mathbf B_0(S))}
 \left\langle\mathbf B'(S)\ \middle|\ \boldsymbol\eta',\mu\right\rangle_d^{\bullet,\,\cC'/(Q'\cup\{p\})}
 \left\langle\mathbf B_0(S)_{\lambda:A\rightsquigarrow B}
 \ \middle|\ \boldsymbol\eta_0,\mu\right\rangle_d^{\bullet,\,\cC_0/(Q_0\cup\{p\})}\Bigg].
\end{aligned}
\]


For
\(\chi_{\log}(\cC,Q)\partial^2/(\partial t^1_q\,\partial t^1_{k-q-2})\),
setting \(A=\tau_q(\omega)\) and \(B=\tau_{k-q-2}(\omega)\), we add \(\{A,B\}\) in the insertions.
Following similar approach for \(-\chi_{\log}(\cC,Q)\,\partial/\partial t^1_k\), we get
\[
\begin{aligned}
&\chi_{\log}(\cC,Q)
 \left\langle\mathbf B\sqcup\{A,B\}
 \ \middle|\ \boldsymbol\eta',\boldsymbol\eta_0\right\rangle_d^{\bullet,\,\cC/Q}
={}\sum_{S\subset\{1,\ldots,u\}}\sum_{\mu\vdash d}\mathfrak z(\mu)\Bigg[
\\[-2mm]
&\quad\chi_{\log}(\cC',Q'\cup\{p\})
 \left\langle\mathbf B'(S)\sqcup\{A,B\}
 \,\middle|\,\boldsymbol\eta',\mu\right\rangle_d^{\bullet,\,\cC'/(Q'\cup\{p\})}
 \left\langle\mathbf B_0(S)\,\middle|\,\boldsymbol\eta_0,\mu\right\rangle_d^{\bullet,\,\cC_0/(Q_0\cup\{p\})}
\\[-1mm]
&\quad+\chi_{\log}(\cC_0,Q_0\cup\{p\})
 \left\langle\mathbf B'(S)\,\middle|\,\boldsymbol\eta',\mu\right\rangle_d^{\bullet,\,\cC'/(Q'\cup\{p\})}
 \left\langle\mathbf B_0(S)\sqcup\{A,B\}
 \,\middle|\,\boldsymbol\eta_0,\mu\right\rangle_d^{\bullet,\,\cC_0/(Q_0\cup\{p\})}\Bigg].
\end{aligned}
\]


For
\(\partial^2/(\partial x^{i,a}_q\,\partial x^{i,a^\vee}_{k-q-1})\), setting
\(A=\tau_q(\phi_{i,a})\) and
\(B=\tau_{k-q-1}(\phi_{i,a^\vee})\),  we add \(\{A,B\}\) in the insertions. Applying \eqref{eq:relative-orbifold-degeneration}, if \(p_i\in\cC'\), then
\[
\begin{aligned}
&\left\langle\mathbf B\sqcup\{A,B\}
 \ \middle|\ \boldsymbol\eta',\boldsymbol\eta_0\right\rangle_d^{\bullet,\,\cC/Q}
\\
={}&\sum_{S\subset\{1,\ldots,u\}}\sum_{\mu\vdash d}\mathfrak z(\mu)
 \left\langle\mathbf B'(S)\sqcup\{A,B\}
 \ \middle|\ \boldsymbol\eta',\mu\right\rangle_d^{\bullet,\,\cC'/(Q'\cup\{p\})}
 \left\langle\mathbf B_0(S)\ \middle|\ \boldsymbol\eta_0,\mu\right\rangle_d^{\bullet,\,\cC_0/(Q_0\cup\{p\})},
\end{aligned}
\]
and if \(p_i\in\cC_0\), then
\[
\begin{aligned}
&\left\langle\mathbf B\sqcup\{A,B\}
 \ \middle|\ \boldsymbol\eta',\boldsymbol\eta_0\right\rangle_d^{\bullet,\,\cC/Q}
\\
={}&\sum_{S\subset\{1,\ldots,u\}}\sum_{\mu\vdash d}\mathfrak z(\mu)
 \left\langle\mathbf B'(S)
 \ \middle|\ \boldsymbol\eta',\mu\right\rangle_d^{\bullet,\,\cC'/(Q'\cup\{p\})}
 \left\langle\mathbf B_0(S)\sqcup\{A,B\}\ \middle|\ \boldsymbol\eta_0,\mu\right\rangle_d^{\bullet,\,\cC_0/(Q_0\cup\{p\})};
\end{aligned}
\]

For \(\chi_{\log}(\cC,Q)(t^0_0)^2\), setting \(A=\tau_0(\id)\), we delete two distinct markings of \(\mathbf B\) with insertion
\(A\).  From \eqref{eq:chi-log-degeneration-in-proof}, we have
\begin{align*}
    &\chi_{\log}(\cC,Q)
 \sum_{\substack{\lambda,\nu\in\Lambda_A(\mathbf B)\\ \lambda\ne\nu}}
 \left\langle\mathbf B\setminus\{\lambda,\nu\}
 \ \middle|\ \boldsymbol\eta',\boldsymbol\eta_0\right\rangle_d^{\bullet,\,\cC/Q}\\
 ={}&\chi_{\log}(\cC',Q'\cup\{p\})
 \sum_{\substack{\lambda,\nu\in\Lambda_A(\mathbf B)\\ \lambda\ne\nu}}
 \left\langle\mathbf B\setminus\{\lambda,\nu\}
 \ \middle|\ \boldsymbol\eta',\boldsymbol\eta_0\right\rangle_d^{\bullet,\,\cC/Q}\\
 &\qquad+\chi_{\log}(\cC_0,Q_0\cup\{p\})
 \sum_{\substack{\lambda,\nu\in\Lambda_A(\mathbf B)\\ \lambda\ne\nu}}
 \left\langle\mathbf B\setminus\{\lambda,\nu\}
 \ \middle|\ \boldsymbol\eta',\boldsymbol\eta_0\right\rangle_d^{\bullet,\,\cC/Q}\\
 ={}&(I)+(II).
\end{align*}
Applying \eqref{eq:relative-orbifold-degeneration}, we get
\begin{align*}
    (I)&=\chi_{\log}(\cC',Q'\cup\{p\})\sum_{\mu\vdash d}\mathfrak z(\mu)\\
    &\cdot\sum_{\substack{S\subset\{1,\ldots,u\}\\\lambda,\nu\in\Lambda_A(\mathbf B'(S)), \lambda\ne\nu}}
 \left\langle\mathbf B'(S)\setminus\{\lambda,\nu\}
 \ \middle|\ \boldsymbol\eta',\mu\right\rangle_d^{\bullet,\,\cC'/(Q'\cup\{p\})}
 \left\langle\mathbf B_0(S)\ \middle|\ \boldsymbol\eta_0,\mu\right\rangle_d^{\bullet,\,\cC_0/(Q_0\cup\{p\})},
\end{align*}
and 
\begin{align*}
    (II)&=\chi_{\log}(\cC_0,Q_0\cup\{p\})\sum_{\mu\vdash d}\mathfrak z(\mu)\\
    &\cdot\sum_{\substack{S\subset\{1,\ldots,u\}\\\lambda,\nu\in\Lambda_A(\mathbf B_0(S)), \lambda\ne\nu}}
 \left\langle\mathbf B'(S)
 \ \middle|\ \boldsymbol\eta',\mu\right\rangle_d^{\bullet,\,\cC'/(Q'\cup\{p\})}
 \left\langle\mathbf B_0(S)\setminus\{\lambda,\nu\}\ \middle|\ \boldsymbol\eta_0,\mu\right\rangle_d^{\bullet,\,\cC_0/(Q_0\cup\{p\})}.
\end{align*}

% For each summand of
% \eqref{eq:chi-log-degeneration-in-proof}, applying the degeneration formula
% with both selected labels on the corresponding component gives
% \[
% \begin{aligned}
% &\chi_{\log}(\cC,Q)
%  \sum_{\substack{\lambda,\nu\in\Lambda_A(\mathbf B)\\ \lambda\ne\nu}}
%  \left\langle\mathbf B\setminus\{\lambda,\nu\}
%  \ \middle|\ \boldsymbol\eta',\boldsymbol\eta_0\right\rangle_d^{\bullet,\,\cC/Q}
% \\
% ={}&\sum_{\mu\vdash d}\mathfrak z(\mu)\Bigg[
% \\[-2mm]
% &\chi_{\log}(\cC',Q'\cup\{p\})
%  \sum_{\substack{S\subset\{1,\ldots,u\}\\\lambda,\nu\in\Lambda_A(\mathbf B'(S))\\ \lambda\ne\nu}}
% \\[-1mm]
% &{}\cdot
%  \left\langle\mathbf B'(S)\setminus\{\lambda,\nu\}
%  \ \middle|\ \boldsymbol\eta',\mu\right\rangle_d^{\bullet,\,\cC'/(Q'\cup\{p\})}
%  \left\langle\mathbf B_0(S)\ \middle|\ \boldsymbol\eta_0,\mu\right\rangle_d^{\bullet,\,\cC_0/(Q_0\cup\{p\})}
% \\[-1mm]
% &+\chi_{\log}(\cC_0,Q_0\cup\{p\})
%  \sum_{\substack{S\subset\{1,\ldots,u\}\\\lambda,\nu\in\Lambda_A(\mathbf B_0(S))\\ \lambda\ne\nu}}
% \\[-1mm]
% &{}\cdot
%  \left\langle\mathbf B'(S)\ \middle|\ \boldsymbol\eta',\mu\right\rangle_d^{\bullet,\,\cC'/(Q'\cup\{p\})}
%  \left\langle\mathbf B_0(S)\setminus\{\lambda,\nu\}
%  \ \middle|\ \boldsymbol\eta_0,\mu\right\rangle_d^{\bullet,\,\cC_0/(Q_0\cup\{p\})}\Bigg].
% \end{aligned}
% \]

For $t^0_0t^1_0$, setting $A=\tau_0(\id)$ and $B=\tau_0(\omega)$, we delete two markings $\lambda\in\Lambda_A(\mathbf B)$ and $\nu\in\Lambda_B(\mathbf B)$ in $\mathbf B$. Applying \eqref{eq:relative-orbifold-degeneration}, if $\nu\in T$, then
\[
\begin{aligned}
 &\left\langle\mathbf B\setminus\{\lambda,\nu\}
 \ \middle|\ \boldsymbol\eta',\boldsymbol\eta_0\right\rangle_d^{\bullet,\,\cC/Q}
\\
={}&\sum_{\mu\vdash d}\mathfrak z(\mu)
\sum_{\substack{S\subset\{1,\ldots,u\}\\\lambda\in\Lambda_A(\mathbf B_0(S))\\\nu\in\Lambda_B(\mathbf B_0(S))}}
 \left\langle\mathbf B'(S)\ \middle|\ \boldsymbol\eta',\mu\right\rangle_d^{\bullet,\,\cC'/(Q'\cup\{p\})}
 \left\langle\mathbf B_0(S)\setminus\{\lambda,\nu\}
 \ \middle|\ \boldsymbol\eta_0,\mu\right\rangle_d^{\bullet,\,\cC_0/(Q_0\cup\{p\})},
\end{aligned}
\]
and if $\nu\in T'$, then
\[
\begin{aligned}
&\left\langle\mathbf B\setminus\{\lambda,\nu\}
 \ \middle|\ \boldsymbol\eta',\boldsymbol\eta_0\right\rangle_d^{\bullet,\,\cC/Q}
\\
={}&\sum_{\mu\vdash d}\mathfrak z(\mu)
 \sum_{\substack{S\subset\{1,\ldots,u\}\\\lambda\in\Lambda_A(\mathbf B'(S))\\\nu\in\Lambda_B(\mathbf B'(S))}}
 \left\langle\mathbf B'(S)\setminus\{\lambda,\nu\}
 \ \middle|\ \boldsymbol\eta',\mu\right\rangle_d^{\bullet,\,\cC'/(Q'\cup\{p\})}
 \left\langle\mathbf B_0(S)\ \middle|\ \boldsymbol\eta_0,\mu\right\rangle_d^{\bullet,\,\cC_0/(Q_0\cup\{p\})}.
\end{aligned}
\]
Observe that, when $\nu\in T$, the condition $\nu\in\Lambda_B(\mathbf B'(S))$ is empty, and when $\nu\in T'$, the condition $\nu\in\Lambda_B(\mathbf B_0(S))$ is empty. Consequently, we obtain
\[
\begin{aligned}
&\sum_{\substack{\lambda\in\Lambda_A(\mathbf B)\\\nu\in\Lambda_B(\mathbf B)}}
 \left\langle\mathbf B\setminus\{\lambda,\nu\}
 \ \middle|\ \boldsymbol\eta',\boldsymbol\eta_0\right\rangle_d^{\bullet,\,\cC/Q}
={}\sum_{\substack{\lambda\in\Lambda_A(\mathbf B)\\\nu\in\Lambda_B(\mathbf B)}}\sum_{\mu\vdash d}\mathfrak z(\mu)\Bigg[
\\[-2mm]
&\quad
 \sum_{\substack{S\subset\{1,\ldots,u\}\\\lambda\in\Lambda_A(\mathbf B'(S))\\\nu\in\Lambda_B(\mathbf B'(S))}}
 \left\langle\mathbf B'(S)\setminus\{\lambda,\nu\}
 \ \middle|\ \boldsymbol\eta',\mu\right\rangle_d^{\bullet,\,\cC'/(Q'\cup\{p\})}
 \left\langle\mathbf B_0(S)\ \middle|\ \boldsymbol\eta_0,\mu\right\rangle_d^{\bullet,\,\cC_0/(Q_0\cup\{p\})}
\\[-1mm]
&\quad+\sum_{\substack{S\subset\{1,\ldots,u\}\\\lambda\in\Lambda_A(\mathbf B_0(S))\\\nu\in\Lambda_B(\mathbf B_0(S))}}
 \left\langle\mathbf B'(S)\ \middle|\ \boldsymbol\eta',\mu\right\rangle_d^{\bullet,\,\cC'/(Q'\cup\{p\})}
 \left\langle\mathbf B_0(S)\setminus\{\lambda,\nu\}
 \ \middle|\ \boldsymbol\eta_0,\mu\right\rangle_d^{\bullet,\,\cC_0/(Q_0\cup\{p\})}\Bigg].
\end{aligned}
\]

For \(x^{i,a}_0x^{i,a^\vee}_0\), setting
\(A=\tau_0(\phi_{i,a})\) and \(B=\tau_0(\phi_{i,a^\vee}))\), 
we delete two markings $\lambda\in\Lambda_A(\mathbf B)$ and $\nu\in\Lambda_B(\mathbf B)$ in $\mathbf B$.
Applying \eqref{eq:relative-orbifold-degeneration}, if $p_i\in\cC'$, then
\[
\begin{aligned}
&\sum_{\substack{\lambda\in\Lambda_A(\mathbf B)\\\nu\in\Lambda_B(\mathbf B)}}\left\langle\mathbf B\setminus\{\lambda,\nu\}
 \ \middle|\ \boldsymbol\eta',\boldsymbol\eta_0\right\rangle_d^{\bullet,\,\cC/Q}
={}\sum_{\substack{\lambda\in\Lambda_A(\mathbf B)\\\nu\in\Lambda_B(\mathbf B)}}
 \sum_{\substack{S\subset\{1,\ldots,u\}\\\lambda\in\Lambda_A(\mathbf B'(S))\\\nu\in\Lambda_B(\mathbf B'(S))}}
 \sum_{\mu\vdash d}\mathfrak z(\mu)
\\[-1mm]
&\hspace*{20mm}\cdot
 \left\langle\mathbf B'(S)\setminus\{\lambda,\nu\}
 \ \middle|\ \boldsymbol\eta',\mu\right\rangle_d^{\bullet,\,\cC'/(Q'\cup\{p\})}
 \left\langle\mathbf B_0(S)\ \middle|\ \boldsymbol\eta_0,\mu\right\rangle_d^{\bullet,\,\cC_0/(Q_0\cup\{p\})},
\end{aligned}
\]
and if $p_i\in\cC_0$, then
\[
\begin{aligned}
&\sum_{\substack{\lambda\in\Lambda_A(\mathbf B)\\\nu\in\Lambda_B(\mathbf B)}}\left\langle\mathbf B\setminus\{\lambda,\nu\}
 \ \middle|\ \boldsymbol\eta',\boldsymbol\eta_0\right\rangle_d^{\bullet,\,\cC/Q}
={}\sum_{\substack{\lambda\in\Lambda_A(\mathbf B)\\\nu\in\Lambda_B(\mathbf B)}}
 \sum_{\substack{S\subset\{1,\ldots,u\}\\\lambda\in\Lambda_A(\mathbf B_0(S))\\\nu\in\Lambda_B(\mathbf B_0(S))}}
 \sum_{\mu\vdash d}\mathfrak z(\mu)
\\[-1mm]
&\hspace*{20mm}\cdot
 \left\langle\mathbf B'(S)
 \ \middle|\ \boldsymbol\eta',\mu\right\rangle_d^{\bullet,\,\cC'/(Q'\cup\{p\})}
 \left\langle\mathbf B_0(S)\setminus\{\lambda,\nu\}\ \middle|\ \boldsymbol\eta_0,\mu\right\rangle_d^{\bullet,\,\cC_0/(Q_0\cup\{p\})},
\end{aligned}
\]

The term \(s^j_0\bar s^j_0\) in \(L_{-1}\) is handled directly at the
level of coefficient extraction.  Since all odd
variables occur on \(\cC'\), it follows that \eqref{eq:relative-orbifold-degeneration} gives
\[
\begin{aligned}
&\mathbf B!\,[\mathbf t_{\mathbf B}]\,s^j_0\bar s^j_0
 Z_d^{\cC/Q}[\boldsymbol\eta]
\\
={}&\sum_{S\subset\{1,\ldots,u\}}\sum_{\mu\vdash d}\mathfrak z(\mu)
 \Bigl(\mathbf B'(S)!\,[\mathbf t_{\mathbf B'(S)}]\,s^j_0\bar s^j_0
 Z_d^{\cC'/(Q'\cup\{p\})}[\boldsymbol\eta',\mu]\Bigr)\cdot
 \left\langle\mathbf B_0(S)\ \middle|\ \boldsymbol\eta_0,\mu\right\rangle_d^{\bullet,\,\cC_0/(Q_0\cup\{p\})}.
\end{aligned}
\]
The relative order of the odd variables is unchanged in
\(\mathbf B'(S)\),
hence the sign is the same on both sides.



For $c_0(\cC/Q)$ in $L_0$, note that
\[
c_0(\cC/Q)
=\sum_{i=1}^{s}\frac{r_i^2-1}{24r_i}
=c_0\bigl(\cC'/(Q'\cup\{p\})\bigr)
+c_0\bigl(\cC_0/(Q_0\cup\{p\})\bigr).
\]
Now  \eqref{eq:relative-orbifold-degeneration} gives
\[
\begin{aligned}
&c_0(\cC/Q)
 \left\langle\mathbf B\ \middle|\ \boldsymbol\eta',\boldsymbol\eta_0
 \right\rangle_d^{\bullet,\,\cC/Q}
={}\sum_{S\subset\{1,\ldots,u\}}\sum_{\mu\vdash d}\mathfrak z(\mu)\Bigg[
\\[-2mm]
&\quad c_0\bigl(\cC'/(Q'\cup\{p\})\bigr)
 \left\langle\mathbf B'(S)\ \middle|\ \boldsymbol\eta',\mu\right\rangle_d^{\bullet,\,\cC'/(Q'\cup\{p\})}
 \left\langle\mathbf B_0(S)\ \middle|\ \boldsymbol\eta_0,\mu\right\rangle_d^{\bullet,\,\cC_0/(Q_0\cup\{p\})}
\\[-1mm]
&\quad+c_0\bigl(\cC_0/(Q_0\cup\{p\})\bigr)
 \left\langle\mathbf B'(S)\ \middle|\ \boldsymbol\eta',\mu\right\rangle_d^{\bullet,\,\cC'/(Q'\cup\{p\})}
 \left\langle\mathbf B_0(S)\ \middle|\ \boldsymbol\eta_0,\mu\right\rangle_d^{\bullet,\,\cC_0/(Q_0\cup\{p\})}\Bigg].
\end{aligned}
\]


Finally, combining the identities above for all operators in \(L_k\) gives
\eqref{eq:virasoro-defect-degeneration}. This completes the proof.
\end{proof}

\end{document}
